\subsection{Liệt dây 7 ngoại biên}

\subsubsection{Chỉ định}
Xoa bóp cho các nguyên nhân sau
\begin{itemize}
    \item Do nhiễm lạnh
    \item Do nhiễm trùng
    \item Do xung huyết (do chấn thương).
\end{itemize}

Xoa bóp chỉ giúp phục hồi nhanh vì vậy cần phối hợp với điều trị
nguyên nhân như thuốc hoặc phẫu thuật.

\subsubsection{Cách làm}

\begin{longtblr}[
    caption = {Liệt dây 7 ngoại biên},
]{
    width = \textwidth,
    colspec = {| X[1,c] |X[1, l]|X[2, l]|X[1, c]|X[7, l]|},
    rowhead = 1,
}
    \hline
    ID & Động tác & Vùng tác động & Số lần & Lưu ý
    \\ \hline
    1 & Xoa & Mặt & 5-10 & {\
        Từ đường trung tâm chia mặt thành hai phần trái và phải.
        Tay trái phần mặt trái.
        Tay phải phần mặt phải.
        Đặt lòng bàn tay xoa theo vòng tròn theo chiều từ trung tâm (khoé mũi)
        sang hai bên (má)
    }
    \\ \hline
    2 & Xát & Mặt & 5-10 & {\
        Từ đường trung tâm chia mặt thành 2 phần trái và phải.
        Tay trái phần mặt trái.
        Tay phải phần mặt phải.
        Thực hiện xát tịnh tiến.
    }
    \\ \hline
    3 & Miết & Mặt & 5-10 & {\
        Mỗi tay một bên mặt, mắt thì miết ngược từ bình minh lên toản trúc
    }
    \\ \hline
    4 & Phân hợp  & Mặt & 5-10 & {
    }
    \\ \hline
    5 & Véo & Lông mày & 5-10 & {\
        Véo chân lông mày
    }
    \\ \hline
    6 & Day & Toả cơ vùng mặt & 5-10 & {
    }
    \\ \hline
    7 & Bóp & Cơ mặt & 5-10 & {
    }
    \\ \hline
    8 & Vờn & Cơ mặt & 5-10 & {
    }
    \\ \hline
    9 & Bấm & Huyệt & 5-10 & {\
        % Nội Đình (ST 44) cùng bên.
        % Hợp Cốc (LI 4)  khác bên.
        % Thừa Tương (CV 24)
        % Địa Thương (ST 4)
        % Nghinh Hương (Ex-HN9)
        % Nhân Trung (GV 26)
        % Giáp Xa (ST 6) %(\giapxa)
        % Hạ Quan (ST 7) %(\haquan)
        % Ế Phong (TH 17)
        % Tình Minh (BL 1) %(\tinhminh)
        % Thừa Khấp (ST 1)
        % Đồng Tử Liêu (GB 1)
        % Ty Trúc Không (TH 23/ SJ 23)
        % Ngư Yêu (Ex-HN4) %(\nguyeu)
        % Toản Trúc (BL 2)
        % Dương Bạch (GB 14)
    }
    \\ \hline
    & & & & {\
        Nội Đình (Cùng bên)
        \includegraphics[width=\linewidth]{ST_44.jpg}
    }
    \\ \hline
    & & & & {\
        Hợp Cốc (Khác bên)
        \includegraphics[width=\linewidth]{LI_4.jpg}
    }
    \\ \hline
    & & & & {\
        Thừa Tương (CV 24)
        \includegraphics[width=\linewidth]{CV 24.jpg}

        Nằm ở chỗ lõm dưới môi dưới và phía trên cằm, tại chính giữa rãnh môi - cằm.
    }
    \\ \hline
    & & & & {\
        Địa Thương (ST 4)
        \includegraphics[width=\linewidth]{ST 4.jpg}

        Huyệt này nằm cách góc miệng (khoé miệng) 0,4 thốn về phía ngoài.
        Huyệt này nằm trên rãnh mũi - má (nếu bình thường không nhìn rõ, mỉm cười sẽ giúp thấy rõ hơn).
    }
    \\ \hline
    & & & & {\
        Nghinh Hương (LI 20)
        \includegraphics[width=\linewidth]{LI 20.jpg}

        Huyệt này nằm ở điểm giữa bờ ngoài của cánh mũi, trong rãnh mũi - má.
    }
    \\ \hline
    & & & & {\
        Nhân Trung (GV 26)
        \includegraphics[width=\linewidth]{GV 26.jpg}

        Nằm ở vị trí 2/3 phía trên bờ môi trên, trên rãnh nhân trung và
        đường giữa mặt.
        Rãnh nhân trung là chỗ lõm nằm trên đường giữa mặt, nằm giữa mũi
        và bờ môi trên.
    }
    \\ \hline
    & & & & {
        Giáp Xa (ST 6)
        \includegraphics[width=\linewidth]{ST 6_1_0.jpg}
        \newline
        \includegraphics[width=\linewidth]{ST 6_2.jpg}

        Huyệt này nằm ở phía trước và phía trên góc hàm khoảng 1 chiều ngang ngón tay, nằm trên cơ cắn.

        Cơ cắn sẽ nổi rõ hơn khi cắn chặt răng (cắn hàm).
    }
    \\ \hline
    & & & & {\
        Hạ Quan (ST 7)
        \includegraphics[width=\linewidth]{ST 7.jpg}
        Huyệt này nằm ở chỗ lõm tại bờ dưới gò má (cung gò má),
        phía trước lồi cầu xương hàm dưới.

        Có thể xác định huyệt bằng cách há miệng và đặt ngón tay
        lên lồi cầu xương hàm dưới, khi ngậm miệng lại ngón tay
        sẽ trượt vào đúng vị trí của huyệt ST-7 (khi châm cứu
        được thực hiện lúc ngậm miệng).
    }
    \\ \hline
    & & & & {\
        Ế Phong (TH 17)
        \includegraphics[width=\linewidth]{TH 17.jpg}
        Huyệt này nằm ở chỗ lõm phía sau dái tai, trung điểm giữa
        mỏm chũm và ngành lên xương hàm dưới, ngay phía trên mỏm ngang
        của đốt sống cổ C1 (đốt sống cổ thứ nhất).
        Dễ xác định huyệt hơn khi ép (gập) dái tai về phía trước.
    }
    \\ \hline
    & & & & {\
        Tình Minh (BL 1)
        \includegraphics[width=\linewidth]{BL 1.jpg}
        Huyệt này nằm ở vị trí phía trên và trong khóe mắt
        trong 0,1 thốn, sát bờ trong xương hốc mắt.
    }
    \\ \hline
    & & & & {\
        Thừa Khấp (ST 1)
        \includegraphics[width=\linewidth]{ST 1.jpg}
        Huyệt này nằm ngay phía dưới đồng tử (khi mắt nhìn thẳng về phía trước),
        ở khoảng giữa nhãn cầu và bờ dưới hốc mắt.
    }
    \\ \hline
    & & & & {\
        Đồng Tử Liêu (GB 1)
        \includegraphics[width=\linewidth]{GB 1.jpg}
        Huyệt này nằm cách khóe mắt ngoài 0,5 thốn về phía bên (ngoài), nằm trong
        chỗ lõm phía ngoài của bờ xương hốc mắt.
    }
    \\ \hline
    & & & & {\
        Ty Trúc Không (TH 23)
        \includegraphics[width=\linewidth]{TH 23.jpg}
        Nằm ở chỗ lõm tại đuôi lông mày, ngay trên bờ trên hốc mắt.
    }
    \\ \hline
    & & & & {\
        Ngư Yêu (Ex-HN4)
        \includegraphics[width=\linewidth]{Ex-HN4.jpg}

        Nằm ở chính giữa lông mày, thẳng hàng với con ngươi (khi mắt nhìn thẳng về phía trước).
    }
    \\ \hline
    & & & & {\
        Toản Trúc (BL 2)
        \includegraphics[width=\linewidth]{BL 2.jpg}
        Huyệt này nằm ở đầu trong lông mày, thẳng ngay phía trên huyệt BL-1
        (Tình Minh), trong chỗ lõm ở khuyết trên hốc mắt.
    }
    \\ \hline
    & & & & {\
        Dương Bạch (GB 14)

        \includegraphics[width=\linewidth]{GB 14.jpg}

        Huyệt này nằm thẳng phía trên con ngươi (khi nhìn thẳng về phía trước)
        và cách điểm giữa lông mày lên phía trên 1 thốn.
    }
    \\ \hline
    10 & Đẩy & Toàn Trúc & &
    \\ \hline
    11 & Miết ngược & & &
    \\ \hline
    12 & Đẩy & Rãnh cười & &
    \\ \hline
\end{longtblr}

\subsubsection{Liệu trình}

\begin{itemize}
    \item Tần suất: ngày 1 lần; mỗi lần gần 30 phút làm đến khi khỏi.
    \item Lấy bên liệt làm chủ nhưng không bỏ qua bên lành.
    \item Giai đoạn đầu: chú trọng bảo vệ - trị mắt, tránh lộ đồng tử kéo dài gây khô/viêm giác mạc - kết mạc.
    \item Khuyên đeo kính râm bảo vệ mắt khi ra ngoài.
\end{itemize}
